\begin{definition}
Với $\lambda>0$ và $y\in\Hil$, điểm gần kề của $y$ theo hàm $F$ được định nghĩa bởi
\begin{equation}
    \prox_{\lambda F}(y)
    :=
    \argmin_{x\in\Hil}
    \left\{
    \Phi_\lambda(x)
    :=
    F(x)
    +
    \frac1{2\lambda}\norm{x-y}^2
    \right\}.
\end{equation}
\end{definition}

\begin{essayonly}
Do $\Phi_\lambda$ là tổng của một hàm lồi và một hàm toàn phương $1/\lambda$-lồi mạnh, bài toán trên luôn có nghiệm duy nhất. Hơn nữa, điều kiện tối ưu bậc một bài toán lồi cho ta đặc trưng quen thuộc
\end{essayonly}
\begin{equation}
  z = \prox_{\lambda F}(y)
  \iff 0 \in \partial \Phi_\lambda(z)
  \iff \frac{y-z}{\lambda} \in \partial F(z),
  \label{eq:opt-cond}
\end{equation}
\begin{essayonly}
Từ đặc trưng trên suy ra $$\prox_{\lambda F}=(I+\lambda\partial F)^{-1},$$ nghĩa là toán tử gần kề chính là toán tử phân giải của dưới vi phân $\partial F$.
\end{essayonly}
Trên thực tế, việc tính $\prox_{\lambda F}(y)$ tự nó có thể khó không kém việc giải bài toán ban đầu, nên ta cần một khái niệm về xấp xỉ của điểm gần kề. Công cụ thường dùng là $\eps$-dưới vi phân: với $\eps\ge 0$,
\begin{equation}
  \partial_\eps F(z) := \big\{ \xi\in\Hil :
     F(x) \ge F(z) + \iprod{x-z}{\xi} - \eps,\ \forall x\in\Hil \big\}.
  \label{eq:eps-subdiff}
\end{equation}
Khi $\eps=0$, ta thu được dưới vi phân thông thường. Ta có một tính chất cơ bản sau đây
\begin{equation}
  0\in\partial_\eps F(z)
  \iff
  F(z)\le\inf F+\eps.
  \label{eq:eps-subopt}
\end{equation}
\begin{essayonly}
Nói cách khác, điều kiện \textit{$0$ là một $\eps$-dưới gradient tại $z$} tương đương với việc \textit{$z$ là một điểm $\eps$-tối ưu} của bài toán.
\begin{proof}
Theo định nghĩa của $\eps$-dưới vi phân,
\[
0\in\partial_\eps F(z)
\iff
F(x)\ge F(z)-\eps,\qquad\forall x\in\Hil.
\]
Lấy infimum theo $x$ suy ra
\[
\inf F\ge F(z)-\eps,
\iff
F(z)\le\inf F+\eps.
\]
Ngược lại, nếu $F(z)\le\inf F+\eps$ thì với mọi $x\in\Hil$,
\[
F(x)\ge\inf F\ge F(z)-\eps
=
F(z)+\langle x-z,0\rangle-\eps,
\]
nên theo định nghĩa của $\partial_\eps F(z)$,
\[
0\in\partial_\eps F(z).
\]
\end{proof}
\end{essayonly}

\begin{essayonly}
Dựa trên khái niệm $\eps$-dưới vi phân, ta sẽ xét ba mô hình xấp xỉ của điểm gần kề. Cả ba đều thuộc lớp sai số tuyệt đối, nghĩa là tham số sai số $\eps$ được quy định độc lập với vị trí của điểm $y$ cũng như độ lớn của điểm gần kề cần tính.
\end{essayonly}

